\uuid{qnYn}
\exo7id{3309}
\titre{exo7 3309}
\auteur{quercia}
\organisation{exo7}
\datecreate{2010-03-09}
\isIndication{false}
\isCorrection{false}
\chapitre{Application linéaire}
\sousChapitre{Définition}
\cpge{mpsi-mp2i.applications-lineaires.b, pcsi.applications-lineaires.b, ptsi.applications-lineaires.b, mpsi-mp2i.espaces-vectoriels.c, pcsi.espaces-vectoriels.c, ptsi.espaces-vectoriels.c}
\module{Algèbre}
\niveau{L1}
\difficulte{}

\contenu{
\texte{
Soit $E$ un espace vectoriel et $f \in \mathcal{L}(E)$ tel que pour tout
$\vec x \in E$, la famille $\bigl(\vec x, f(\vec x)\bigr)$ est liée.
}
\begin{enumerate}
    \item \question{Montrer que si $\vec x \ne \vec 0$, il existe un unique scalaire
    $\lambda_{\vec x}$ tel que $f(\vec x) = \lambda_{\vec x}\vec x$.}
    \item \question{Comparer $\lambda_{\vec x}$ et $\lambda_{\vec y}$ lorsque
    $(\vec x, \vec y)$ est libre.}
    \item \question{Montrer que $f$ est une homothétie.}
\end{enumerate}
}
